\documentclass[a4paper,12pt]{article}
\usepackage[utf8]{inputenc}
\usepackage[italian]{babel}
\usepackage{listings}
\usepackage{xcolor}
\usepackage{float}
\usepackage{placeins}
\usepackage{amssymb}
\setlength{\emergencystretch}{3em}

\lstset{
  language=C,
  backgroundcolor=\color{lightgray!20},
  frame=single,
  breaklines=true,
  numbers=left,
  numberstyle=\tiny,
  stepnumber=1,
  numbersep=5pt,
  captionpos=b,
  basicstyle=\ttfamily\small,
  keywordstyle=\color{blue}\bfseries,
  commentstyle=\color{green!60!black},
  stringstyle=\color{red},
  showstringspaces=false
}

\title{LabReSiD26 -- Hands-On 7}
\author{Davide Ferrara -- 518629}
\date{}

\begin{document}

\maketitle

% ---------------------------------------------------------------------------
\section{Esercizio 1: Osserva il comportamento base}
% ---------------------------------------------------------------------------

\subsection*{Domanda 1 -- Quanti receiver ricevono il messaggio e perché TCP unicast non lo permetterebbe}

Eseguendo tre receiver all'interno di Tmux in tre terminali diversi e inviando un
messaggio tramite sender, si può notare che tutti i receiver ricevono il messaggio.
Con TCP unicast questo non sarebbe possibile perché con esso si stabilisce una
connessione 1-a-1 e non 1-a-N come nel caso del multicast: il sender dovrebbe aprire
N connessioni separate, una per ogni receiver, e inviare il messaggio N volte.
Con multicast è sufficiente una sola \texttt{sendto()} — il kernel e la rete si
occupano della distribuzione.

\subsection*{Domanda 2 -- MAC di destinazione catturato con tcpdump per 239.0.0.1}

Output catturato con \texttt{sudo tcpdump -i enp8s0 udp port 5000 -e}:

\begin{lstlisting}
13:21:03.385919 74:56:3c:c7:5e:5d (oui Unknown) > 01:00:5e:00:00:01 (oui Unknown),
ethertype IPv4 (0x0800), length 60:
dave.lan.48598 > 239.0.0.1.commplex-main: UDP, length 5
\end{lstlisting}

Il MAC di destinazione \texttt{01:00:5e:00:00:01} corrisponde con quello atteso.
Applicando la formula \texttt{01:00:5e:XX:XX:XX}: gli ultimi 23 bit di \texttt{239.0.0.1}
(\texttt{EF|00|00|01}) sono \texttt{0x000001}, quindi il MAC risultante è
\texttt{01:00:5e:00:00:01}.

\subsection*{Domanda 3 -- Nuovo MAC di destinazione con gruppo 239.1.2.3}

Modificando \texttt{MCAST\_GROUP} in \texttt{239.1.2.3} in sender.c e receiver.c,
ricompilando e catturando con tcpdump:

\begin{lstlisting}
13:31:45.237811 74:56:3c:c7:5e:5d (oui Unknown) > 01:00:5e:01:02:03 (oui Unknown),
ethertype IPv4 (0x0800), length 60:
dave.lan.50275 > 239.1.2.3.commplex-main: UDP, length 5
\end{lstlisting}

Il nuovo MAC di destinazione è \texttt{01:00:5e:01:02:03}. Gli ultimi 23 bit di
\texttt{239.1.2.3} corrispondono agli utlimi 3 ottetti di \texttt{MCAST\_GROUP}.

% ---------------------------------------------------------------------------
\section{Esercizio 2: TTL e loopback}
% ---------------------------------------------------------------------------

\subsection*{Domanda 1 -- Cosa succede con \texttt{ttl = 0} e perché}

Con \texttt{ttl = 0}, \texttt{sendto()} ritorna successo e il sender stampa
\textit{Inviato (x byte)...}, ma il pacchetto non raggiunge nessuna interfaccia fisica:
il kernel lo scarta al livello IP prima che venga consegnato al driver di rete e
\texttt{tcpdump} sull'interfaccia fisica non cattura nulla.

La differenza tra i due valori è definita qui sotto:

\begin{itemize}
  \item \textbf{TTL = 0} — scope host-local: il pacchetto non esce dalla macchina.
        Il kernel lo sopprime prima che raggiunga qualsiasi interfaccia di rete.
        Nessun altro host lo riceve.
  \item \textbf{TTL = 1} — scope link-local: il pacchetto esce sull'interfaccia
        di rete ma non viene instradato oltre il primo hop. Raggiunge tutti gli
        host sulla stessa LAN iscritti al gruppo, ma nessun router lo forwarda.
\end{itemize}

\newpage
\subsection*{Domanda 2 -- Effetto di \texttt{IP\_MULTICAST\_LOOP = 0} sulla stessa macchina e su macchina diversa}

Con \texttt{loop = 0} il kernel non consegna una copia del pacchetto UDP ai socket
locali iscritti al gruppo. Effetto nei due scenari:

\begin{itemize}
  \item \textbf{Stessa macchina:} il receiver non riceve nulla. \texttt{IP\_MULTICAST\_LOOP}
        sopprime la copia locale prima che venga messa nella coda del socket.
  \item \textbf{Macchina diversa nella stessa LAN:} il receiver riceve normalmente.
        L'opzione non influenza il pacchetto che esce sull'interfaccia fisica: questo
        viene inviato sulla rete con TTL=1 e raggiunge tutti gli host della LAN
        iscritti al gruppo.
\end{itemize}

\subsection*{Domanda 3 -- Differenza tra \texttt{IP\_MULTICAST\_TTL} e
\texttt{IP\_MULTICAST\_LOOP}}

Dal manuale:

\begin{verbatim}
IP_MULTICAST_TTL: Set or read the time-to-live value of
outgoing multicast packets for this socket. The default is
1 which means that multicast packets don't leave the local
network unless the user program explicitly requests it.

IP_MULTICAST_LOOP: Set or read a boolean integer argument
that determines whether sent multicast packets should be
looped back to the local sockets.
\end{verbatim}

\textbf{IP\_MULTICAST\_TTL} controlla quanti router il pacchetto può attraversare
sulla rete. Con TTL=1 (default) resta nella LAN locale. Con TTL=0 il kernel sopprime
il pacchetto silenziosamente prima di consegnarlo al driver di rete: \texttt{sendto()} ritorna
successo ma nessun frame esce sull'interfaccia fisica.

\textbf{IP\_MULTICAST\_LOOP} controlla se il kernel consegna una copia del pacchetto
ai socket locali prima che esca sulla rete. Non influenza l'uscita del pacchetto
sull'interfaccia fisica: con \texttt{loop = 0} il pacchetto viene comunque inviato
sulla rete, ma nessun socket locale lo riceve.

% ---------------------------------------------------------------------------
\section{Esercizio 3: Costruisci una chat multicast}
% ---------------------------------------------------------------------------

\subsection*{Passo A -- Un solo socket per inviare e ricevere}

\texttt{chat.c} è costruito a partire dalla struttura di \texttt{receiver.c}
(\texttt{bind()} + \texttt{IP\_ADD\_MEMBERSHIP}), viene utilizzato
lo stesso scoket sia per recvfrom() che per sendto()

\subsection*{Passo B -- Il problema del blocco}

Mettendo \texttt{fgets()} e \texttt{recvfrom()} in sequenza nel loop:

\begin{lstlisting}
while (1) {
    fgets(buf, sizeof(buf), stdin);    /* blocca finche' l'utente scrive */
    sendto(...);
    recvfrom(sockfd, buf, ...);        /* blocca finche' arriva un pacchetto */
    printf(...);
}
\end{lstlisting}

Il programma si blocca su \texttt{fgets()} finché l'utente non scrive, e non può
ricevere messaggi dalla rete nel frattempo. Simmetricamente, mentre aspetta su
\texttt{recvfrom()} non può leggere da tastiera. Le due operazioni bloccanti in
sequenza rendono il programma inutilizzabile come chat.

\subsection*{Passo C -- \texttt{select()} come soluzione}

Con \texttt{select()} su \texttt{STDIN\_FILENO} e \texttt{sockfd} ad ogni
iterazione il kernel segnala quale dei due fd è pronto e solo quello 
viene servito:

\begin{lstlisting}
while (1) {
    fd_set readfds;
    FD_ZERO(&readfds);
    FD_SET(STDIN_FILENO, &readfds);
    FD_SET(sockfd, &readfds);
    select(sockfd + 1, &readfds, NULL, NULL, NULL);

    if (FD_ISSET(STDIN_FILENO, &readfds)) { /* tastiera pronta: invia */ }
    if (FD_ISSET(sockfd, &readfds))        { /* rete pronta: ricevi  */ }
}
\end{lstlisting}

Invio e ricezione funzionano contemporaneamente senza bloccarsi.

\subsection*{Passo D -- Il proprio messaggio torna indietro}

Per filtrare i propri messaggi è stato adottato
il metodo a): il payload viene prefissato con \texttt{[username]} e alla ricezione
si controlla se il messaggio inizia con il proprio prefisso:

\begin{lstlisting}
/* invio */
snprintf(buf, sizeof(buf), "[%s] %s", username, input);
sendto(sockfd, buf, strlen(buf), 0, (struct sockaddr *)&local, sizeof(local));

/* ricezione */
char own_prefix[BUFLEN];
snprintf(own_prefix, sizeof(own_prefix), "[%s]", username);
if (strncmp(buf, own_prefix, strlen(own_prefix)) == 0)
    continue;
\end{lstlisting}

La differenza rispetto a \texttt{IP\_MULTICAST\_LOOP = 0}, in questo caso si
filtra in userspace dopo la ricezione, quindi i propri messaggi arrivano comunque al
socket ma vengono scartati.

\end{document}
